\subsection{ODE Foundations and Flows}
% Existence/uniqueness, flows, phase portraits, invariant sets.

Ordinary differential equations can often be rewritten as fixed-point problems on a complete function space.

\begin{definition} \label{def:scalar-initial-value-problem} \hlt{Scalar Initial-Value Problem}\\
Let $O\subseteq\R^2$ be open, $(x_0,y_0)\in O$, and $g:O\to\R$. The scalar initial-value problem determined by $(O,g,x_0,y_0)$ is $y'=g(x,y)$ with $y(x_0)=y_0$. For an interval $J\subseteq\R$, a function $f:J\to\R$ solves this problem on $J$ iff
\begin{enumerate}[label=(\roman*)]
\setlength{\itemsep}{0pt}
\item $J$ is open, $x_0\in J$, and $f$ is differentiable.
\item $\Gamma_f\defeq\{(x,f(x)):x\in J\}\subseteq O$.
\item $f'(x)=g(x,f(x))$ for every $x\in J$, and $f(x_0)=y_0$.
\end{enumerate}
\hlt{Structural Commentary}: The graph condition makes the right-hand side defined along $f$; the initial condition selects the integral curve through $(x_0,y_0)$.
\end{definition}

\begin{proposition} \label{prop:ode-integral-equation-equivalence} \hlt{Differential and Integral Forms}\\
Let $O\subseteq\R^2$ be open, $(x_0,y_0)\in O$, $g\in C(O,\R)$, $J\subseteq\R$ be open interval with $x_0\in J$. For $f\in C(J,\R)$ with $\Gamma_f\defeq\{(x,f(x)):x\in J\}\subseteq O$, $f$ solves $y'=g(x,y)$, $y(x_0)=y_0$, on $J$ iff $f(x)=y_0+\int_{x_0}^{x}g(t,f(t))\,dt$ $\forall x\in J$.
\begin{proof}
\begin{enumerate}[label=\roman*.]
\setlength{\itemsep}{0pt}
\item[$(\Rightarrow)$] If $f'(x)=g(x,f(x))$ on $J$ and $f(x_0)=y_0$, then by Fundamental Theorem of Calculus (Theorem \ref{thm:fundamental-theorem-calculus-riemann}),\\
$f(x)-y_0=\int_{x_0}^{x}f'(t)\,dt=\int_{x_0}^{x}g(t,f(t))\,dt$ $\forall x\in J$.
\item[$(\Leftarrow)$] If $f(x)=y_0+\int_{x_0}^{x}g(t,f(t))\,dt$ for every $x\in J$, then $g(\,\cdot\,,f(\,\cdot\,))$ is continuous on $J$.\\
By Theorem \ref{thm:integral-function-continuity-and-derivative}, $f'(x)=g(x,f(x))$ on $J$. Evaluating at $x=x_0$ gives $f(x_0)=y_0$.
\end{enumerate}
\hlt{Structural Commentary}: Differentiation and integration are inverse along the continuous graph.\\
\hlt{Key Insight}: The IVP is equivalent to its Volterra integral equation.
\end{proof}
\end{proposition}

\begin{theorem} \label{thm:picard-local-existence-uniqueness} \hlt{Picard Local Existence Theorem}\\
Let $O\subseteq\R^2$ be open, $(x_0,y_0)\in O$, and $g\in C(O,\R)$.\\
Assume $\abs{g(x,y_1)-g(x,y_2)}\leq M\abs{y_1-y_2}$ for some $M>0$ whenever $(x,y_1),(x,y_2)\in O$.\\
Then $\exists a,b,\ell>0$, $Q\defeq[x_0-a,x_0+a]\times[y_0-b,y_0+b]\subseteq O$, and $J_\ell\defeq(x_0-\ell,x_0+\ell)$ such that the IVP $y'=g(x,y)$, $y(x_0)=y_0$, has a unique differentiable solution $f:J_\ell\to\R$ with $\Gamma_f\subseteq Q$.
\begin{proof}
\begin{enumerate}[label=(\roman*)]
\setlength{\itemsep}{0pt}
\item \hlt{Local box}: Choose $a,b>0$ with $Q\defeq[x_0-a,x_0+a]\times[y_0-b,y_0+b]\subseteq O$, choose $K<\infty$ with $\abs{g}\leq K$ on $Q$, and choose $\ell>0$ with $\ell\leq a$, $\ell K\leq b$, and $\ell M<1$. Set $I_\ell\defeq[x_0-\ell,x_0+\ell]$.
\item \hlt{Complete space}: Let $\mathcal{X}_\ell\defeq\{u\in C(I_\ell,\R):\abs{u-y_0}_\infty\leq b\}$. This is a complete closed subset of $C(I_\ell,\R)$ by Theorem \ref{thm:complete-supremum-metric-space} and Proposition \ref{prop:basic-completeness-consequences}. Define $(Tu)(x)\defeq y_0+\int_{x_0}^{x}g(t,u(t))\,dt$.
\item \hlt{Contraction.} For $u\in\mathcal{X}_\ell$, $\abs{(Tu)(x)-y_0}\leq \abs{x-x_0}K\leq b$, so $T(\mathcal{X}_\ell)\subseteq\mathcal{X}_\ell$.\\
For $u,v\in\mathcal{X}_\ell$, $\norm{Tu-Tv}_{\infty}\leq \ell M\norm{u-v}_{\infty}$, so $T$ is a contraction.
\item \hlt{Fixed point}: Theorem \ref{thm:contraction-principle} gives a unique fixed point $u_*\in\mathcal{X}_\ell$.\\
With $f\defeq u_*|_{J_\ell}$, Proposition \ref{prop:ode-integral-equation-equivalence} gives $f'=g(\,\cdot\,,f)$ and $f(x_0)=y_0$, while $\Gamma_f\subseteq Q$.
\item \hlt{Uniqueness}: If $h:J_\ell\to\R$ is another solution with $\Gamma_h\subseteq Q$, then $\forall 0<r<\ell$ the same operator on $I_r\defeq[x_0-r,x_0+r]$ is a contraction, and both $h|_{I_r}$ and $u_*|_{I_r}$ are fixed points.\\
Hence $h=f$ on every $I_r$, and therefore on $J_\ell$.
\end{enumerate}
\hlt{Structural Commentary}: Box control plus short-time contraction.\\
\hlt{Key Insight}: Picard iteration is Banach iteration.
\end{proof}
\end{theorem}

\begin{figure}[H]
  \centering
  \begin{tikzpicture}[
  >=Stealth,
  x=0.92cm,
  y=0.92cm,
  every node/.style={font=\scriptsize},
  axis/.style={->, black!65, line width=0.55pt},
  rect/.style={draw=blue!55!black, fill=blue!5, fill opacity=0.55, line width=0.65pt},
  strip/.style={draw=orange!80!black, fill=orange!8, fill opacity=0.55, line width=0.6pt},
  graph/.style={blue!70!black, line width=0.8pt},
  iter/.style={orange!85!black, line width=0.6pt, dashed},
  guide/.style={black!35, dashed, line width=0.45pt},
  point/.style={circle, fill=orange!85!black, inner sep=1.3pt},
  note/.style={fill=white, inner sep=1.1pt, text=black!72, align=center}
]
  \path[use as bounding box] (-2.55,-1.70) rectangle (3.40,1.78);

  \draw[rect] (-2.00,-1.12) rectangle (2.00,1.12);
  \draw[strip] (-1.25,-1.12) rectangle (1.25,1.12);

  \draw[axis] (-2.45,0) -- (2.65,0) node[right] {$x$};
  \draw[axis] (0,-1.50) -- (0,1.60) node[above] {$y$};

  \node[note, anchor=south west] at (2.02,1.10)
    {\textcolor{blue!55!black}{$Q$}};
  \node[note] at (0,-1.42) {\textcolor{orange!80!black}{$I_\ell$}};

  \draw[guide] (-1.25,-1.12) -- (-1.25,-1.34);
  \draw[guide] (1.25,-1.12) -- (1.25,-1.34);
  \node[note, anchor=north] at (-1.25,-1.35) {$x_0-\ell$};
  \node[note, anchor=north] at (1.25,-1.35) {$x_0+\ell$};

  \draw[iter]
    plot[smooth, tension=0.75] coordinates {
      (-1.25,0.32) (-0.72,0.18) (-0.25,0.08)
      (0,0) (0.48,0.14) (0.92,0.27) (1.25,0.36)
    };
  \draw[graph]
    plot[smooth, tension=0.75] coordinates {
      (-1.25,-0.62) (-0.74,-0.32) (-0.30,-0.16)
      (0,0) (0.44,0.28) (0.92,0.47) (1.25,0.55)
    };

  \node[point] at (0,0) {};
  \node[note, anchor=south west] at (0.08,0.54) {$(x_0,y_0)$};
  \node[note, anchor=south west] at (-1.18,0.36)
    {\textcolor{orange!85!black}{$f_n$}};
  \node[note, anchor=west] at (1.38,0.62)
    {\textcolor{blue!70!black}{$f_{n+1}=Tf_n$}};
\end{tikzpicture}

  \caption{Picard setup: choose $\ell$ with $\ell K\leq b$ so $T(\mathcal X_\ell)\subseteq\mathcal X_\ell$ and graphs stay in $Q$.}
  \label{fig:picard-rectangle-integral-operator}
\end{figure}

\begin{method} \label{meth:picard-iteration-local-ode} \hlt{Picard Iteration}\\
Let $Q\defeq[x_0-a,x_0+a]\times[y_0-b,y_0+b]\subseteq O$, assume $\abs{g(x,y)}\leq K$ for $(x,y)\in Q$,\\
$\abs{g(x,y_1)-g(x,y_2)}\leq M\abs{y_1-y_2}$ whenever $(x,y_1),(x,y_2)\in Q$. Let $\ell>0$ with $\ell\leq a$, $\ell K\leq b$, and $\ell M<1$.\\
Set $I_\ell\defeq[x_0-\ell,x_0+\ell]$, $J_\ell\defeq(x_0-\ell,x_0+\ell)$, $\mathcal{X}_\ell\defeq\{u\in C(I_\ell,\R):\abs{u-y_0}_\infty\leq b\}$,\\
$(Tu)(x)\defeq y_0+\int_{x_0}^{x}g(t,u(t))\,dt$. Then
\begin{enumerate}[label=(\roman*)]
\setlength{\itemsep}{0pt}
\item Initialize $f_0(x)\equiv y_0$ on $I_\ell$.
\item Iterate $f_{n+1}\defeq Tf_n$, i.e. $f_{n+1}(x)=y_0+\int_{x_0}^{x}g(t,f_n(t))\,dt$ for $x\in I_\ell$.
\item Let $f_\infty\defeq\lim_{n\to\infty}f_n$ uniformly on $I_\ell$, and output $f\defeq f_\infty|_{J_\ell}$.
\end{enumerate}
Since $T$ is a contraction on $\mathcal{X}_\ell$, the limit exists; the output is the unique local solution with graph in $Q$.
\end{method}
